\section{Introduction}

Clinical text is scarce, but structured medical knowledge is not. France maintains medical ontologies that span tens of thousands of codes (CIM-10 for diagnoses, CCAM for procedures, ATC for medications), and PubMed Central contains millions of clinical case descriptions embedded in scientific articles. Both are publicly available, both are richly structured, and neither can be used directly to train a language model. The question is whether we can convert them into synthetic training data that is closer to what clinical applications actually need.

We already saw one form of generated data in \Cref{chap:ontobook}, where medical ontologies became synthetic textbooks for pretraining. This chapter generates data for the extraction task itself, in two ways, and neither touches a real patient record. The first annotates public biomedical text with a large language model, producing the labelled examples that train the extraction model of \Cref{chap:architectures}. The second rewrites public case descriptions into the style of hospital notes, to approximate a register that no public corpus offers. Both take a public source, hand it to a language model under a strict prompt, and get back training material that real clinical data could not supply.


\section{Annotating Biomedical Text for Extraction}
\label{sec:gliner-data}

\begin{figure}[t]
\centering
\providecommand{\ftitle}[1]{{\normalsize\bfseries\sffamily\color{ThesisInk}#1}}
\providecommand{\etype}[1]{{\color{ThesisPrimaryDark}#1}}
\resizebox{\textwidth}{!}{%
\begin{tikzpicture}[
  font=\small\sffamily,
  outbox/.style={draw, rounded corners=4pt, line width=0.8pt,
                 text width=7.4cm, inner sep=12pt, align=left, font=\small},
]
% Input text block (left)
\node[draw=ThesisNeutral, fill=ThesisPaper, rounded corners=4pt, line width=0.8pt,
      text width=3.7cm, inner sep=11pt, align=justify, font=\small,
      text=ThesisInk] (input) at (0,0)
  {\itshape Mrs Diallo, 64, diffuse large B-cell lymphoma diagnosed 12/03/2019.
   R-CHOP started 02/04/2019, 6 cycles. Relapse 15/01/2021. Death 03/02/2021.
   LDH 480 IU/L at diagnosis.};

% Three output blocks (right), stacked with generous spacing
\node[outbox, draw=ThesisPrimaryDark, fill=ThesisPrimary!10,
      anchor=west] (ent) at (5.6,4.2)
  {\ftitle{Entities}\par\vspace{5pt}
   {\ttfamily\etype{Disease}\,$\rightarrow$\,diffuse large B-cell lymphoma\par
    \etype{Drug}\,$\rightarrow$\,R-CHOP\quad
    \etype{Examination}\,$\rightarrow$\,LDH\par
    \etype{date}\,$\rightarrow$\,03/02/2021\quad
    \etype{measure}\,$\rightarrow$\,480 IU/L}};

\node[outbox, draw=ThesisSecondaryDark, fill=ThesisSecondary!12,
      anchor=west] (rel) at (5.6,0)
  {\ftitle{Relations}\par\vspace{5pt}
   {\ttfamily treats:\, R-CHOP\,$\rightarrow$\,lymphoma}};

\node[outbox, draw=ThesisTertiaryDark, fill=ThesisTertiary!14,
      anchor=west] (str) at (5.6,-4)
  {\ftitle{Structure \& classification}\par\vspace{5pt}
   {\ttfamily timeline \{diagnosis date, treatment start,\par
    \hphantom{timeline \{}relapse, death\}\par
    ICD-10 chapter: tumours\par
    document type: follow-up note}};

% Fan-out arrows through an LLM label
\coordinate (hub) at (4.2,0);
\draw[-{Latex[length=2.6mm]}, line width=1.1pt, ThesisInk!55]
   (input.east) -- node[above, font=\small\bfseries, color=ThesisInk]{LLM} (hub);
\draw[-{Latex[length=2.6mm]}, line width=1pt, ThesisInk!45] (hub) to[out=55,in=180] (ent.west);
\draw[-{Latex[length=2.6mm]}, line width=1pt, ThesisInk!45] (hub) -- (rel.west);
\draw[-{Latex[length=2.6mm]}, line width=1pt, ThesisInk!45] (hub) to[out=-55,in=180] (str.west);
\end{tikzpicture}%
}
\caption{A large language model turns one passage into typed entities, relations
between them, and a document structure with its classifications. The full
annotation prompts are given in \Cref{app:gliner-prompts}.}
\label{fig:annotation-example}
\end{figure}

Training an extractor by fine-tuning needs labelled examples, and as this part's introduction noted, those are out of reach for French clinical text. We take a different route. Instead of paying clinicians to annotate real notes, we annotate public biomedical text automatically with a large language model, and use the result to train a smaller extraction model. This is a form of distillation: a large general model that is too heavy to run at scale produces the labels, and a small model learns to reproduce them. Everything we annotate is public text, so none of the resulting data carries the restrictions of a patient record.

The annotator is \texttt{Qwen3-235B-A22B-Instruct}, a 235-billion-parameter instruction-tuned model~\citep{yang2025qwen3}, run with FP8 quantisation on four NVIDIA H100 GPUs through vLLM~\citep{kwon2023vllm}. For each document, the model returns a single JSON object whose fields are fixed in advance by a schema, so every annotation has the same shape and parses without error. We turn off the model's step-by-step reasoning mode and sample at a low temperature of $0.6$, which keeps the output close to the schema. The same model later serves as the judge that scores our extractors (\Cref{chap:evaluation}), so we check there that it does not simply favour its own annotations.

The schema asks for more than a list of entities. For each document the model first lists candidate spans, meaning every noun phrase that could be a medical entity, and then builds four kinds of annotation on top of them. It marks \emph{entities} with open-vocabulary types at four levels of granularity, from a single reusable word such as \textit{Maladie} (disease) up to a full compositional phrase, and the same span may be typed at several levels at once. It assigns two or three \emph{classifications} to the document, such as its specialty or its document type. It fills zero to two \emph{structures}, small field-value schemas such as a dosage or a patient profile. And it records \emph{relations} between entities, such as a drug that treats a disease. The model also returns a short list of \emph{hard negatives}: plausible types that are absent from the text, which teach the extractor what not to find. \Cref{fig:annotation-example} shows one annotated example.

We then ground every annotation against the source text. A mention is kept only when it is an exact substring of the document, which removes spans the model paraphrased or invented. We also drop a type when it is merely a lexical variant of its own mention, measured by edit distance, so the model cannot take the shortcut of repeating a word as its own type. The head and tail of every relation must likewise be exact substrings. Documents where no annotation survives are not discarded: a document with text but no entity is kept as a negative example, which is what teaches the extractor to stay silent on non-medical text.

The text we annotate is drawn from three kinds of public source, listed in \Cref{tab:gliner-sources}. Most of it is French biomedical literature from the corpus of \Cref{chap:mcbio}, kept above a quality threshold. A large share is case-centred: articles that contain patient cases, and French translations of PubMed Central case reports. A smaller part comes from passages generated from the CIM-10, CCAM, and ATC nomenclatures, which bring in coding terminology. Finally, non-medical and low-quality documents are added on purpose as negatives. The final dataset holds about $94{,}000$ annotated documents, with $662{,}000$ entities, $1.98$~million mentions, and $592{,}000$ relations; roughly a quarter of the documents are negatives.

\begin{table}[t]
  \centering
  \caption{Public sources sampled for annotation by the language model to train
  the extraction model of \Cref{chap:architectures}. Of the $118{,}000$ sampled
  documents, about $94{,}000$ yield a valid annotation after grounding and
  filtering.}
  \label{tab:gliner-sources}
  \small
  \setlength{\tabcolsep}{12pt}
  \renewcommand{\arraystretch}{1.25}
  \begin{tabular}{@{}ll S[table-format=6.0, group-separator={,}, group-minimum-digits=4]@{}}
    \toprule
    Source & What it provides & {Sampled} \\
    \midrule
    \tablegroup{3}{Biomedical literature (\Cref{chap:mcbio})}
    HAL                 & French biomedical articles             & 26000 \\
    ISTEX               & French biomedical articles             & 26000 \\
    \tablegroup{3}{Patient cases}
    Articles with cases & passages describing real patient cases  & 18000 \\
    PMC-Patients-FR     & translated PubMed Central case reports  & 14000 \\
    \tablegroup{3}{Ontology terminology}
    CIM-10 / CCAM / ATC & passages from the coding nomenclatures  & 9000 \\
    \tablegroup{3}{Negatives}
    Non-medical / low quality & documents with no entity to find  & 25000 \\
    \midrule
    Documents sampled    &                                        & 118000 \\
    Valid annotations    & after grounding and filtering          & 94000 \\
    \bottomrule
  \end{tabular}
\end{table}



\section{Synthetic Clinical Reports}
\label{sec:synthetic-reports}

Public case descriptions still read like published cases, not like the terse notes a hospital produces. The second use of generation closes part of that gap. We rewrite a public case into the register of a hospital note, and a deterministic degradation step adds the abbreviations, missing accents, and dictation slips that make real notes hard to read, so the extractor meets that register during training as well as on real notes.


\section{Synthesizing Longitudinal Patient Dossiers}
\label{sec:lymphoma-dossiers}

A single note captures one moment. A patient's record is a story that unfolds over years and across many documents, and the clinical questions we care about live in that story rather than in any one note. This thesis is part of OncoLab, a project that annotates the dossiers of lymphoma patients, all of them relapsed, against a form of about a hundred variables covering demographics, disease history, treatment lines, imaging, and outcome. Real dossiers cannot be used to develop and test that annotation, for the privacy reasons that run through this thesis, so the project needs a synthetic but realistic set of French dossiers to build the form and train automatic extraction before any real record is touched.

A first attempt seeded each synthetic dossier from a single short case description, and it did not work. The reason is worth stating, because it shaped everything that followed. The seed text was about a thousand characters and described one moment, while a dossier spans years and a dozen documents. Asked to bridge that distance, the model had to invent almost everything, and it did so by imitation, producing dossiers that looked plausible but did not hold together. Clinicians who reviewed the output found the revealing gaps: there were almost no follow-up reports and almost no relapse consultations, although every patient in the target cohort has relapsed; cases with both a first and a second treatment line were missing, so adverse events beyond the first cycle could not be annotated; and the multidisciplinary-meeting reports had the wrong shape.

The fix was not to add a rule for each defect. Patching the prompt with one constraint per problem makes the model produce clone-like patients and loses the variety that made it useful. The better question is what lets the model do well on its own, and the answer was a better source and a decomposed task. We seed instead from real longitudinal case reports, translated into French, in which relapse, second-line treatment, adverse events, and death are already present in the text. From the $167{,}034$ public cases we keep the richest, giving more weight to those that mention a relapse and a second treatment line, down to about $2{,}800$ adult cases.

The generation then splits into stages and respects time. The model first writes a coherent patient timeline as structured data, and a short reflexive pass repairs any inconsistency in that timeline rather than rejecting it. Only then does it write the documents, one at a time and in chronological order: each document is generated seeing the earlier ones but never the later ones, so the dossier accumulates the way a real record does and no document can quote a result from its own future. A final pass extracts the structured fields. Because a field such as a treatment date takes several values over a dossier, each field is stored as a list of values, each tied to the document it came from, which keeps the longitudinal structure and lets every value be traced back to its source.

The redesign produced about three thousand synthetic dossiers, around eight documents each, with a valid timeline and a complete annotation in almost every case. Second-line treatment now appears in about a third of the dossiers, against roughly a fifth before, which is the gap the clinicians had flagged. The work is still in progress: the most recent run covers $2{,}050$ of its $2{,}774$ planned cases, and the automatic quality check measures format coherence rather than clinical accuracy, so it should not be read as a correctness score. A held-out part of this set becomes the hardest evaluation in the thesis, the one we use in \Cref{chap:evaluation}.


\section*{Conclusion}

Across these uses the recipe is the same: take a public source, constrain a language model with a strict prompt, and keep only what holds together. The data this produces is real and large, but it settles nothing on its own. It only matters if a model can turn it into clinical extraction under a label set that is open and keeps changing, the setting that defeats a fixed classifier.

\vspace{2em}
The data now exists. Whether a model can use it well is the question of the next chapter, \Cref{chap:architectures}, which builds the extractor trained on the material generated here.
